% !TEX root = ../../Main.tex
\section{Experimental Procedure: Training, Selection and Deployment}
\label{Section:ExperimentalProcedure}

An agent design does not produce a result on its own. Training is stochastic, so the same configuration run twice gives two different policies. A procedure is needed to decide which of them is reported, and why. This section describes that procedure, the criteria used at each stage, and the two deployment decisions that sit between the learned policy and the orders actually sent.

\subsection{Training Procedure: Seeds, Episodes and Data Splits}

The eleven-year window is split once into three consecutive periods. The window is written as $[T_0, T_3)$, where $T_0$ is the first of January 2015 and $T_3$ is the first of January 2026:

\begin{equation}
\underbrace{[T_0, T_1)}_{\text{training, 9 years}} \quad
\underbrace{[T_1, T_2)}_{\text{validation, 1 year}} \quad
\underbrace{[T_2, T_3)}_{\text{testing, 1 year}}
\label{Equation:Split}
\end{equation}

Weights are updated on the training period only. The validation period decides which state of the network is kept. The testing period is used once, after every other decision has been made. Within this split, a single training run works as follows. A seed fixes the network initialisation, the stream of exploration noise and the order in which the replay buffer is sampled, so it determines the whole run. The agent then runs $30$ episodes. An episode is one pass over the training period. Weights are not reset between episodes. The run is a single sequence of $30$ passes, and the network at the start of episode $n+1$ is the network at the end of episode $n$. After each episode the agent is evaluated once over the validation period, with learning turned off. This evaluation is reduced to a single number, the Calmar ratio, which is the annualised return divided by the maximum drawdown. Let $f_n$ be this value and let $g_n$ be an eligibility test defined below. The state kept after episode $n$ is

\begin{equation}
\theta^{\star} \leftarrow \theta_n \quad \text{if} \quad (g_n \wedge \neg g^{\star}) \;\vee\; \left(g_n = g^{\star} \wedge f_n > f^{\star}\right)
\label{Equation:KeepBest}
\end{equation}

which prefers an eligible episode to an ineligible one, whatever the scores. Otherwise it keeps the higher score within the same eligibility class. The saved state is a copy. The run continues training from $\theta_n$, not from $\theta^{\star}$, so a later episode never restarts from an earlier checkpoint. The eligibility test $g_n$ rejects policies that are degenerate. It does not reject a policy only because it performs poorly. An episode passes only if it opened at least ten positions, spent at least three hundred bars on each side of the market, and held its weaker side for at least $30\%$ of the time it held its stronger one. A policy that never trades, or that is long for the whole period, can reach a good Calmar ratio on a trending window without having learned anything about when to hold a position. The test removes such a policy before the scores are compared. Between $128$ and $256$ seeds were run for each currency pair, and the runs differed only in the seed. \Cref{Figure:TrainingPipeline} shows how the three stages fit together.

\begin{figure}[htb!]
\centering
\begin{tikzpicture}[
  node distance=5mm,
  stage/.style={draw, rounded corners=2pt, align=center, font=\small, text width=32mm, minimum height=8mm, inner sep=3pt},
  small/.style={stage, font=\scriptsize, text width=30mm},
  note/.style={font=\scriptsize\itshape, align=left},
  flow/.style={-stealth, thick}
]
\node[stage] (seed) {Seed $s$\\\scriptsize init, noise, sampling};
\node[small, right=10mm of seed] (episode) {Episode $n$ of 30\\one pass over training};
\node[small, right=10mm of episode] (valid) {Validation pass\\learning disabled};
\node[small, below=9mm of valid] (keep) {Calmar $f_n$ + eligibility\\keep-best checkpoint};
\node[small, below=9mm of episode] (test) {Test pass\\once, after 30 episodes};
\node[stage, below=9mm of seed] (pool) {Pool of 128--256\\trained seeds};
\node[stage, below=9mm of pool] (gates) {Gates\\reject degenerate};
\node[small, right=10mm of gates] (sweep) {Risk sweep\\$\times 1.0$ to $\times 3.0$};
\node[small, right=10mm of sweep] (score) {Composite score\\percentile ranks};
\node[stage, below=9mm of gates] (champion) {Published model};

\draw[flow] (seed) -- (episode);
\draw[flow] (episode) -- (valid);
\draw[flow] (valid) -- (keep);
\draw[flow] (keep) -- node[note, above, midway] {} (test);
\draw[flow] (keep.west) to[out=180, in=0] node[note, above, midway] {\scriptsize next episode} (episode.south east);
\draw[flow] (test) -- (pool);
\draw[flow] (pool) -- (gates);
\draw[flow] (gates) -- (sweep);
\draw[flow] (sweep) -- (score);
\draw[flow] (score.south) to[out=270, in=0] (champion.east);
\end{tikzpicture}
\caption{The three stages of the procedure. The Calmar ratio chooses among the episodes of one seed. The composite score of \Cref{Section:ExperimentalProcedure} chooses among seeds.}
\label{Figure:TrainingPipeline}
\end{figure}

One property of this setup must be stated clearly, because it limits what the results in \Cref{Chapter:ExperimentalResults} can claim. The split of \Cref{Equation:Split} applies to the training of a single seed. The choice of which seed to publish, described next, is based on performance over the whole eleven-year window. The performance figures reported in this work are measured over that same window. The testing period is therefore held out from the training of each candidate and from the choice among its episodes. It is not held out from the choice among candidates, so the headline returns are not out-of-sample estimates. This is the weakest point of the method, and it is not hidden anywhere in the rest of this thesis.

\subsection{Model Selection: Gates and Composite Ranking}

Once the seeds for a pair have been trained, each one is evaluated over the full window with the cost model of \Cref{Tables:CostModel}, and one is chosen. The choice is made in two steps, first rejection and then ranking, because they answer two different questions. The first question is whether a policy is admissible at all. The second is which of the admissible policies is best.

Four gates reject a candidate directly. First, its return must be positive. Second, its weaker direction must make up at least $10\%$ of its exposure. This removes policies that are long or short the whole time. Third, its mean holding period must be at least twenty-four bars. This removes policies whose signal is noise. Fourth, its regime score must be higher than the score of the same policy with its position sequence rotated in time. The regime score is the fraction of market movement during which the policy held the correct position. The last gate needs some explanation, because a fixed threshold was tried first and does not work. A rotated policy holds the same positions for the same durations, but at the wrong moments. Its score therefore measures what the policy's autocorrelation and holding pattern would earn by luck alone. Across the seven pairs these null scores range from $48$ to $70$. A fixed threshold of, for example, $55$ would be below chance on one pair and too strict on another. The null score is therefore computed for each model and not assumed.

Candidates that pass the gates are ranked by a composite of percentile ranks within their own pool. For a metric $m$ and a pool of $K$ candidates, the rank of candidate $i$ is

\begin{equation}
R_m(i) = 100 \cdot \frac{\left|\left\{ j : v_m(j) < v_m(i) \right\}\right|}{K - 1}
\label{Equation:PercentileRank}
\end{equation}

where $v_m$ is negated for metrics that are better when smaller, such as drawdown. The metrics are split into three groups, listed in \Cref{Tables:Composite}. The score is the unweighted mean of the three group means:

\begin{equation}
S(i) = \frac{1}{3} \sum_{G \in \{\text{money}, \text{safety}, \text{behaviour}\}} \frac{1}{|G|} \sum_{m \in G} R_m(i)
\label{Equation:Composite}
\end{equation}

% !TEX root = ../Main.tex
\begin{table}[htb!]
\centering
\footnotesize
\caption{Parts of the composite selection score of \Cref{Equation:Composite}. The three groups have equal weight, so return makes up one ninth of the total and does not dominate it. The ratios are those defined in \Cref{Section:PerformanceMeasurement}.}
\label{Tables:Composite}
\begin{tabularx}{\textwidth}{@{}lX@{}}
\toprule
\textbf{Group} & \textbf{Metrics, each used as a percentile rank within the candidate pool} \\
\midrule
Money & Annualised return, Calmar ratio, Sterling ratio \\
\addlinespace
Safety & Sharpe ratio, Sortino ratio, inverse maximum drawdown, inverse downside volatility \\
\addlinespace
Behaviour & Two-sidedness (the weaker exposure side as a fraction of half), regime edge over the model's own rotation null, mean holding period, longest holding period \\
\bottomrule
\end{tabularx}
\end{table}

Two design decisions are built into \Cref{Equation:Composite}. First, ranks are relative to the pool and not absolute. No threshold can therefore exclude a candidate that is simply better than all the others in its pool. Second, the three groups have equal weight, so a policy cannot win on return alone. Money, safety and behaviour each count for a third. A candidate with the highest return in its pool, but with no two-sidedness and no regime edge over its rotation null, scores poorly on two groups out of three. Each candidate that passes the gates is also evaluated at several position sizes, from the reference risk of \Cref{Equation:PositionSizing} up to three times that risk. Scaling the risk fraction changes only the size of the positions the policy takes. It does not change the learned policy. The level chosen for each pair is reported with the results.

\subsection{Deployment: Decision Frequency and Rebalance Threshold}

Two decisions sit between the learned policy and the orders that are actually sent. Together they explain more of the delivered performance than any hyperparameter.

The first is the decision frequency, which is how often the agent may act. The agent observes every hourly bar, but it may act only once per calendar day. A key is built from the timestamp of each bar, and a new action is taken only when that key changes:

\begin{equation}
\kappa(t) = (\text{year}(t), \text{month}(t), \text{day}(t))
\label{Equation:DecisionKey}
\end{equation}

The key is built from the calendar and not from a count of bars. This means that the gaps and short sessions described in \Cref{Section:TradingFramework} do not affect it. A week with a holiday still produces one decision per trading day, and the decisions stay aligned with the trading days. The second decision is a rebalance threshold. A continuous action changes slightly at every decision. Acting on every such change would send a stream of tiny orders, and each of them would pay the spread and a commission. An order is therefore sent only when the change in the target position is large enough to be worth its cost:

\begin{equation}
\left| V_{target} - V_{current} \right| \;\geq\; \max\!\left(V_{min},\; \eta \left| V_{ref} \right|\right), \qquad \eta = 0.20
\label{Equation:Hysteresis}
\end{equation}

where $V_{min}$ is the broker's minimum volume. The threshold is a fraction of the reference volume, not of the position currently held, so it does not shrink to zero as a position is reduced. \Cref{Tables:Deployment} isolates the effect of each decision, measured on one model with nothing retrained. The comparison is between four configurations of the same weights, so the differences are attributable to the deployment layer alone.

% !TEX root = ../Main.tex
\begin{table}[htb!]
\centering
\footnotesize
\caption{Effect of the two deployment decisions, measured on one EUR/USD model over the full window with nothing retrained. All four rows use the same weights with different execution settings, so the differences are attributable to the deployment layer alone. This model comes from before the seven-pair campaign and is used here only to isolate the effect of each decision.}
\label{Tables:Deployment}
\begin{tabular}{@{}llrrrr@{}}
\toprule
\textbf{Frequency} & \textbf{Threshold} & \textbf{Commission} & \textbf{Return} & \textbf{Sharpe} & \textbf{Max drawdown} \\
\midrule
Daily & --- & none & $+42.48\%$ & $0.319$ & $24.4\%$ \\
Hourly & --- & none & $-38.57\%$ & $-0.308$ & $50.5\%$ \\
Hourly & --- & 3.5 points & $-64.70\%$ & $-0.732$ & $70.5\%$ \\
Daily & $0.20$ & 3.5 points & $+34.02\%$ & $0.275$ & $23.5\%$ \\
\bottomrule
\end{tabular}
\end{table}

Two things follow from that table. First, moving from hourly to daily decisions adds about eighty percentage points and turns a losing configuration into a profitable one. The reason is that hourly decisions act on noise that the eleven-year horizon of the discount factor was never meant to capture. Second, the rebalance threshold recovers about twelve percentage points under realistic costs, and at the same time it reduces drawdown instead of increasing it. This is unusual. Commission alone costs over twenty percentage points without the filter and about six with it. The continuous action space makes the strategy expressive, and this filter keeps its trading costs low enough.

\subsection{Worked Example: One Position From Decision to Settlement}

The earlier parts of this chapter define each component separately. This section follows one position through all of them, using a trade taken by the delivered EUR/USD model. This allows the arithmetic to be checked instead of taken on trust.

At 02:00 on 30 January 2019 the daily decision key of \Cref{Equation:DecisionKey} changed, and the agent was asked for an action for the first time that day. It received the 38-value observation of \Cref{Section:AgentDesign}, built from the bar that had just closed, and it returned a positive action. At that moment the reference volume of \Cref{Equation:PositionSizing} risked $1\%$ of the balance over a stop of $1.5$ average true ranges. The action scaled this volume, and the result was rounded to the broker's volume step. This gave a target of $12{,}000$ units of the base currency. The change from the position held at that time was larger than the threshold of \Cref{Equation:Hysteresis}, so an order was sent. The tick recorded at that instant quoted $1.14392$ ask and $1.14387$ bid, a spread of five points. The buy was filled at the ask. Twenty-four hours later the key changed again, and the agent's action changed sign. The position was closed at the bid quoted at that time, $1.14931$. The profit follows from those two prices and the volume. It is expressed in the account currency at the conversion rate recorded on the closing tick:

\begin{equation}
12{,}000 \times \left(1.14931 - 1.14392\right) = 64.68 \text{ USD} \;\;\longrightarrow\;\; 56.27 \text{ EUR}
\label{Equation:WorkedGross}
\end{equation}

The spread is already included in that figure. It was paid at entry, by filling at the ask and not at the mid, and again at exit, by filling at the bid. This is why the round trip crosses the spread exactly once. The commission is charged separately at both ends:

\begin{equation}
2 \times 3.5 \times 10^{-5} \times 12{,}000 = 0.84 \text{ USD} \;\;\longrightarrow\;\; 0.72 \text{ EUR}
\label{Equation:WorkedCommission}
\end{equation}

The position was held overnight, but the swap is zero because the account is swap-free. The position therefore settled at $56.27 - 0.72 = 55.55$ EUR, which is the figure in the exported trade record. One trade says little on its own. The delivered EUR/USD model took $1{,}777$ trades over the eleven years. Across all of them, the same accounting gives a gross profit of $8{,}733.62$ EUR, a commission charge of $1{,}669.50$ EUR and a swap charge of exactly zero, for a net of $7{,}064.12$ EUR. Commission alone took $19.1\%$ of the gross profit. This figure does not include the spread, which cannot be separated from the gross figure because it is part of every fill. A result reported without these charges would describe a different strategy.